\documentclass[11pt]{article}

\usepackage{amsmath,amssymb}
\usepackage{xurl}
\usepackage[hidelinks]{hyperref}
\setlength{\emergencystretch}{2em}

% Shared commands for notes in docs/paper-gaps/.
% Notation follows the LDT paper (references/ldt-paper/).

\newcommand{\Lean}{\textsc{Lean}}
\newcommand{\leanid}[1]{\textsf{#1}}
\newcommand{\ghissue}[1]{\href{https://github.com/LionSR/MIPStarRE/issues/#1}{GitHub issue \##1}}
\newcommand{\ghpr}[1]{\href{https://github.com/LionSR/MIPStarRE/pull/#1}{GitHub PR \##1}}

% --- Paper-compatible notation ---
\newcommand{\F}{\mathbb{F}}
\newcommand{\N}{\mathbb{N}}
\newcommand{\C}{\mathbb{C}}
\newcommand{\tr}{\operatorname{tr}}
\newcommand{\ot}{\otimes}
\newcommand{\E}{\mathop{\bf E\/}}
\newcommand{\bx}{\mathbf{x}}
\newcommand{\by}{\mathbf{y}}
\newcommand{\bu}{\mathbf{u}}
\newcommand{\bv}{\mathbf{v}}

% Approximation relation (paper uses \approx_\eps and \simeq_\zeta)
\newcommand{\approxeps}{\approx_{\varepsilon}}
\newcommand{\simgeq}{\simeq_{\zeta}}

% Polynomial spaces (paper uses \polyfunc, \polysub, \polymeas)
\newcommand{\polyfunc}[3]{\operatorname{Poly}(#1,#2,#3)}
\newcommand{\polysub}[3]{\operatorname{PolySub}(#1,#2,#3)}
\newcommand{\polymeas}[3]{\operatorname{PolyMeas}(#1,#2,#3)}


\title{Issue 1100: Endpoint Branches in the Proof of
\texorpdfstring{$\path{lem:projective-non-measurement}$}{lem:projective-non-measurement}}
\author{}
\date{2026-05-03}

\begin{document}
\maketitle

\section{At a glance}

\begin{itemize}
  \item \textbf{Difficulty.} Low.  The paper proof treats the nontrivial
  small-error regime explicitly, while the endpoint \(\zeta=0\) and the
  large-error regime require separate formal branches.
  \item \textbf{Estimated weight.} A local spectral endpoint argument together
  with a trivial large-error construction; both are already represented in the
  current Lean development.
  \item \textbf{Mathlib/project split.} Mostly project-local.  Mathlib supplies
  the spectral calculus and matrix-order facts, while the project supplies the
  measurement, state-dependent distance, and spectral-truncation interfaces.
  \item \textbf{Key mathematical inputs.} The source almost-projectivity estimate
  \path{eq:A-looks-projective}, the small-error truncation theorem, the
  \(\zeta=0\) spectral projection argument, and the large-error fallback.
\end{itemize}

\section{Key theorem forms}

The formalization separates the local theorem used in the paper from the
unconditional theorem exposed to later Lean code.
\begin{enumerate}
  \item \textbf{Small-error source branch.}  Under the hypotheses of
  \path{eq:A-looks-projective} and the standing assumption
  $0 < \zeta \le 1/4$, construct projectors $R_a$ satisfying the
  state-dependent-distance estimate and
  $\sum_a R_a \leq (1 + 2\sqrt{\zeta}) I$.
  \item \textbf{Endpoint branch.}  When $\zeta = 0$, replace the positive
  truncation threshold by the spectral projector onto the $1$-eigenspace and
  prove the exact state-dependent-distance conclusion.
  \item \textbf{Unconditional packaging theorem.}  Combine the small-error
  branch, the endpoint branch, and the trivial large-error branch into the
  standalone projective-non-measurement witness consumed by the
  orthonormalization construction.\footnote{The relevant declarations are
  \leanid{projectiveNonMeasurement\_of\_sourceAlmostProjective},
  \leanid{projectiveNonMeasurement\_of\_sourceAlmostProjective\_full}, and
  \leanid{spectralTruncationStatement\_of\_sourceAlmostProjective} in
  \doclink{MIPStarRE/LDT/MakingMeasurementsProjective/SpectralTruncation.lean}.}
\end{enumerate}

\section{The cited statement}

This note concerns the lemma ``Rounding to projectors'' in
\cite[Section~5]{Ji2020LowIndividualDegree}.\footnote{Tracked in
\ghissue{1100}. In the TeX source consulted here, the statement and proof are in
\path{references/ldt-paper/orthonormalization.tex}, lines~414--531.}
The cited lemma states that, for a measurement $A = \{A_a\}$,
there exists a family of projectors $\{R_a\}$ such that
\[
  A_a \ot I \approx_{2\sqrt{\zeta}} R_a \ot I
\]
and, writing $R = \sum_a R_a$,
\[
  R \leq (1 + 2\sqrt{\zeta}) I.
\]
In the surrounding proof, the parameter $\zeta$ is the error term already fixed
in the orthonormalization-main-lemma argument.

Just before this lemma, the paper derives
\[
  \sum_a \bra{\psi} (A_a - A_a^2) \ot I \ket{\psi} \le 2\zeta.
  \tag{\path{eq:A-looks-projective}}
\]
This is the only input used in the proof of the rounding lemma: once
\path{eq:A-looks-projective} has been obtained, the auxiliary measurement $B$ no
longer appears.

\section{The large-error branch}

The paper does handle the large-$\zeta$ region, but it does so one level above
the rounding lemma.  At the start of the proof of the orthonormalization-main
lemma it says that the argument is trivial when $\zeta > 1/4$, and therefore
continues under the standing assumption
\[
  \zeta \le 1/4.
  \tag{\path{eq:assumption-on-zeta}}
\]
This reduction appears in
\path{references/ldt-paper/orthonormalization.tex}, lines~375--379, before the
derivation of \path{eq:A-looks-projective} and before the proof of
\path{lem:projective-non-measurement}.  Thus the cited argument for the rounding
lemma is intended to be read only in the small-error regime.

For the paper as written, this is enough: the lemma is invoked only inside the
proof which has already restricted to $\zeta \le 1/4$.  A standalone formal
declaration, however, cannot rely on that surrounding reduction unless it is
built into the declaration itself.  If one wants a public theorem whose
conclusion is the rounding witness for all $\zeta \ge 0$, then the trivial
large-error branch must be supplied explicitly in the formal development even
though the paper leaves it implicit at the level of
\path{lem:projective-non-measurement}.

\section{The endpoint \texorpdfstring{$\zeta = 0$}{zeta = 0}}

The proof of the rounding lemma introduces an auxiliary threshold parameter
$\delta$ satisfying
\[
  0 < \delta \le 1/2.
  \tag{\path{eq:bound-on-delta}}
\]
The positivity of $\delta$ is used throughout the truncation argument, because
the scalar inequality
\[
  (x - \mathsf{trunc}_\delta(x))^2 \le \frac{1}{\delta}(x - x^2)
\]
and its operator consequence
\[
  (A_a - R_a)^2 \le \frac{1}{\delta}(A_a - A_a^2)
\]
both divide by $\delta$.  The same factor $1/\delta$ is then carried into the
state-dependent-distance estimate in lines~487--493 of the source file.

At the end of the proof, the paper sets
\[
  \delta = \sqrt{\zeta}.
\]
This substitution appears in
\path{references/ldt-paper/orthonormalization.tex}, lines~528--529.  Together
with \path{eq:bound-on-delta}, it covers exactly the regime
\[
  0 < \zeta \le 1/4.
\]
When $\zeta = 0$, the statement of the lemma still makes sense, but the printed
proof no longer applies literally, because it would require $\delta = 0$ while
simultaneously using $1/\delta$.

This is therefore a proof-level endpoint omission rather than a theorem-level
counterexample.  The mathematical content of the $\zeta = 0$ case is clear.  If
\[
  \sum_a \bra{\psi} (A_a - A_a^2) \ot I \ket{\psi} = 0,
\]
then each summand is a nonnegative expectation and hence vanishes.  Writing
$A_a = \sum_i \lambda_{a,i} \ket{u_{a,i}}\bra{u_{a,i}}$ with
$0 \le \lambda_{a,i} \le 1$, one may define $R_a$ to be the spectral projector
onto the $1$-eigenspace of $A_a$.  Then $R_a \le A_a$, so
\[
  \sum_a R_a \le \sum_a A_a = I,
\]
which is stronger than the bound claimed in the lemma.  Moreover the vanishing
of
\(
\bra{\psi} (A_a - A_a^2) \ot I \ket{\psi}
\)
forces the state to have no weight on eigenspaces with eigenvalue strictly
between $0$ and $1$, so
\[
  A_a \ot I \approx_0 R_a \ot I.
\]
The endpoint is therefore valid, but it requires a separate spectral argument.

\section{Relation with the formal development}

The current \Lean{} declarations make this distinction explicit.\footnote{The
constructive small-error branch is formalized by
\leanid{projectiveNonMeasurement\_of\_sourceAlmostProjective} and
\leanid{projectiveNonMeasurement\_of\_almostProjMeasStatement} in
\doclink{MIPStarRE/LDT/MakingMeasurementsProjective/SpectralTruncation.lean}.
The unconditional spectral-truncation producer
\leanid{projectiveNonMeasurement\_of\_sourceAlmostProjective\_full} combines
that branch with the exact $\zeta = 0$ endpoint and the trivial large-error
case, and
\leanid{spectralTruncationStatement\_of\_sourceAlmostProjective} records the
result directly as the spectral-truncation statement used by the later
Section~5 construction.}
Thus the present formal development has discharged the unconditional
spectral-truncation packaging.  The small-error proof still corresponds to the
part of the cited argument in which one may set $\delta = \sqrt{\zeta}$ and
satisfy \path{eq:bound-on-delta}; the zero and large-error regimes are supplied
by separate formal branches.  The remaining Section~5 work lies in the later
locality-preserving repair and completion arguments, not in the endpoint
packaging isolated in this note.

In other words, the formal development has not discovered a contradiction in the
paper statement.  It has isolated the exact places where the cited proof uses
positivity of $\delta$, and therefore where the proof must be supplemented if
one wants a theorem stated without side conditions.

\section{Conclusion}

The paper does not need a separate large-$\zeta$ calculation inside the proof of
\path{lem:projective-non-measurement}, because that regime has already been
removed in the surrounding proof of the orthonormalization-main lemma.  The
formal development must nevertheless add such a branch if it exposes the
rounding lemma as an unconditional standalone theorem.

The more delicate omission is the endpoint $\zeta = 0$.  The statement of the
rounding lemma remains meaningful there, but the printed proof covers only
$0 < \zeta \le 1/4$, because it introduces a positive parameter $\delta$ and
later sets $\delta = \sqrt{\zeta}$.  Thus the cited source leaves the endpoint
argument implicit.  A faithful unconditional formalization should keep the
paper's nontrivial branch for $0 < \zeta \le 1/4$, and add separate proofs for
the large-error region and for the exact endpoint $\zeta = 0$.

\bibliographystyle{alpha}
\bibliography{references}

\end{document}
